\chapter{Μεθοδολογία}
\label{chap:methodology}

\section{Σκοπός και ερευνητικός σχεδιασμός του κεφαλαίου}

Το κεφάλαιο αυτό περιγράφει τη μεθοδολογική λογική της εργασίας, δηλαδή τον
τρόπο με τον οποίο το πρόβλημα του κατακερματισμού πληροφορίας μετατρέπεται σε
ελεγχόμενη ροή συλλογής, εξαγωγής και αξιοποίησης δεδομένων. Δεν αναλύονται
εδώ οι λεπτομέρειες κώδικα, τα αρχεία, οι βιβλιοθήκες ή οι επιμέρους
παράμετροι εκτέλεσης. Οι λεπτομέρειες αυτές ανήκουν στο κεφάλαιο της
υλοποίησης.

Η βασική μεθοδολογική αρχή είναι ότι το γλωσσικό μοντέλο δεν λειτουργεί ως
αυτόνομη πηγή γνώσης. Χρησιμοποιείται ως εργαλείο επεξεργασίας συγκεκριμένου
περιεχομένου που έχει συλλεχθεί από γνωστή πηγή. Αυτή η αρχή αντιστοιχεί σε
εξαγωγή τεκμηριωμένη στην πηγή, όπου το μοντέλο περιορίζεται από το αρχικό
περιεχόμενο και όχι από τη γενική γνώση του. Κάθε αποτέλεσμα οφείλει να
συνδέεται με URL, να ελέγχεται από διαχειριστή και να αποθηκεύεται σε κοινή
δομή ώστε να είναι αξιοποιήσιμο από τη διαδικτυακή εφαρμογή.

\section{Πεδίο εφαρμογής και απαιτήσεις}

\subsection{Πεδίο εφαρμογής και μονάδα ανάλυσης}

Ως μονάδα ανάλυσης ορίζεται ένα χρηματοδοτικό πρόγραμμα ή προϊόν που σχετίζεται
με ενεργειακή αναβάθμιση κατοικιών στην Ελλάδα. Στο πεδίο αυτό περιλαμβάνονται
δημόσια προγράμματα επιχορήγησης, πράσινα ή επιδοτούμενα τραπεζικά προϊόντα
και εργαλεία που συνδέονται με οικιακές παρεμβάσεις εξοικονόμησης ενέργειας ή
ανανεώσιμων πηγών. Δεν εξετάζονται γενικά στεγαστικά δάνεια, απλές
ανακαινίσεις χωρίς ενεργειακό κριτήριο, προγράμματα ηλεκτροκίνησης ή
χρηματοδοτήσεις για καθαρά επιχειρησιακές/βιομηχανικές εφαρμογές.

Η μεθοδολογία δεν αξιολογεί αν ένα πρόγραμμα είναι οικονομικά βέλτιστο για
συγκεκριμένο πολίτη, καθώς αυτό θα απαιτούσε στοιχεία για την κατοικία, το
εισόδημα, τις παρεμβάσεις και το κόστος υλοποίησης. Περιορίζεται στο προηγούμενο
στάδιο: συλλογή, τυποποίηση και παρουσίαση διαθέσιμης πληροφορίας, ώστε αυτή να
μπορεί να αξιοποιηθεί από χρήστη, μηχανικό ή διαχειριστή.

\subsection{Απαιτήσεις της μεθοδολογίας}

Η μεθοδολογία σχεδιάζεται με βάση πέντε απαιτήσεις: επικαιρότητα,
ιχνηλασιμότητα, τυποποίηση, επιλεκτικότητα και ανθρώπινη εποπτεία. Οι απαιτήσεις
αυτές προκύπτουν από τη φύση του προβλήματος: οι πηγές αλλάζουν, οι όροι δεν
εμφανίζονται σε ενιαία μορφή και η πληροφορία μπορεί να έχει οικονομικές ή
διοικητικές συνέπειες.

Οι απαιτήσεις και η πρακτική τους σημασία συνοψίζονται στον
Πίνακα~\ref{tab:methodology-requirements}.

\begin{table}[!htbp]
    \centering
    \caption[Βασικές απαιτήσεις της μεθοδολογίας]{Βασικές απαιτήσεις που καθοδηγούν τη μεθοδολογία.}
    \label{tab:methodology-requirements}
    \begin{tabularx}{\textwidth}{L{3.1cm}Y Y}
        \toprule
        \textbf{Απαίτηση} & \textbf{Μεθοδολογική σημασία} & \textbf{Πρακτική συνέπεια} \\
        \midrule
        Επικαιρότητα & Τα προγράμματα και οι όροι μεταβάλλονται. & Η διαδικασία πρέπει να μπορεί να επαναληφθεί για νέες ή ενημερωμένες πηγές. \\
        Ιχνηλασιμότητα & Κάθε πεδίο πρέπει να ελέγχεται στην πηγή του. & Διατηρούνται URL, χρόνος συλλογής και στοιχεία προέλευσης. \\
        Τυποποίηση & Οι πηγές έχουν διαφορετική μορφή και ορολογία. & Η πληροφορία αποθηκεύεται σε κοινό JSON σχήμα. \\
        Επιλεκτικότητα & Πολλές πηγές είναι οριακά ή καθόλου σχετικές. & Προηγείται απόφαση συνάφειας πριν από την ένταξη στο σύνολο εγγραφών. \\
        Ανθρώπινη εποπτεία & Η αυτόματη εξαγωγή μπορεί να είναι αβέβαιη. & Οι εγγραφές ελέγχονται πριν εμφανιστούν στον τελικό χρήστη. \\
        \bottomrule
    \end{tabularx}
\end{table}

\section{Προτεινόμενη μεθοδολογική ροή}

\subsection{Πηγές και αρχιτεκτονική}

Η μεθοδολογία ξεκινά από πηγές που έχουν θεσμική ή επιχειρησιακή ευθύνη για την
πληροφορία που δημοσιεύουν. Για δημόσια προγράμματα προτιμώνται επίσημες
πλατφόρμες, υπουργικές σελίδες και οδηγοί εφαρμογής. Για τραπεζικά προϊόντα
προτιμώνται οι σελίδες των ίδιων των τραπεζών. Δευτερεύουσες πηγές μπορούν να
βοηθήσουν στον εντοπισμό υλικού, αλλά δεν αποτελούν τελική τεκμηρίωση χωρίς
διασταύρωση.

Η διάκριση των πηγών δεν είναι μόνο οργανωτική. Οι δημόσιες και κρατικές πηγές
τεκμηριώνουν τους όρους των προγραμμάτων, οι τραπεζικές πηγές περιγράφουν τα
χρηματοδοτικά χαρακτηριστικά προϊόντων, ενώ οι θεσμικές ή δευτερεύουσες πηγές
χρησιμοποιούνται κυρίως για πλαίσιο, ανακάλυψη υλικού ή διασταύρωση. Με αυτόν
τον τρόπο, η μεθοδολογία διαχωρίζει την πηγή που τεκμηριώνει μια εγγραφή από
την πηγή που απλώς βοηθά στον εντοπισμό της.

Αρχιτεκτονικά, η προσέγγιση οργανώνεται σε τρία επίπεδα. Το πρώτο επίπεδο
αφορά την προέλευση της πληροφορίας, δηλαδή την επιλογή και αποθήκευση
συγκεκριμένων πηγών. Το δεύτερο αφορά τον μετασχηματισμό της πληροφορίας:
καθαρισμός περιεχομένου, απόφαση συνάφειας, εξαγωγή πεδίων και έλεγχος
ποιότητας. Το τρίτο αφορά την αξιοποίηση, όπου οι ελεγμένες εγγραφές
χρησιμοποιούνται από τη δομημένη αποθήκευση JSON και τη διαδικτυακή εφαρμογή.

Ο διαχωρισμός αυτός είναι απαραίτητος, επειδή αποτρέπει τη σύγχυση ανάμεσα
στην πρωτογενή πηγή, στο αποτέλεσμα του LLM και στην πληροφορία που τελικά
βλέπει ο χρήστης. Έτσι, κάθε στάδιο έχει σαφή ρόλο και μπορεί να ελεγχθεί
ανεξάρτητα.

Η συνολική αρχιτεκτονική της προσέγγισης παρουσιάζεται στο
Σχήμα~\ref{fig:methodology-architecture}, το οποίο συνοψίζει τα στάδια από την
επιλογή πηγών μέχρι την αξιοποίηση των εγκεκριμένων δεδομένων στην εφαρμογή.

\begin{figure}[!htbp]
    \centering
    \begin{adjustbox}{max width=\textwidth,center}
        \begin{tikzpicture}[
    main/.style={
        rectangle,
        rounded corners=3pt,
        draw=blue!55,
        fill=blue!6,
        align=left,
        text width=9.15cm,
        minimum height=0.72cm,
        inner xsep=7pt,
        font=\sffamily\small
    },
    check/.style={
        rectangle,
        rounded corners=3pt,
        draw=blue!55,
        fill=blue!6,
        align=left,
        text width=9.15cm,
        minimum height=0.72cm,
        inner xsep=7pt,
        font=\sffamily\small
    },
    arrow/.style={-{Latex[length=2mm]}, thick, draw=gray!65}
]
\node[main] (sources) {1. Επιλογή επίσημων και τραπεζικών πηγών};
\node[main, below=0.28cm of sources] (snapshot) {2. Συλλογή, καθαρισμός και αποθήκευση στιγμιότυπου περιεχομένου};
\node[main, below=0.28cm of snapshot] (relevance) {3. Έλεγχος συνάφειας με LLM και απόρριψη μη σχετικών πηγών};
\node[main, below=0.28cm of relevance] (extract) {4. Εξαγωγή πεδίων σε κοινό JSON σχήμα};
\node[check, below=0.28cm of extract] (quality) {5. Δομικός και σημασιολογικός έλεγχος με URL, χρόνο και αρχικό κείμενο};
\node[check, below=0.28cm of quality] (admin) {6. Ανθρώπινη εποπτεία, διόρθωση, έγκριση και ανάδραση στη ροή};
\node[main, below=0.28cm of admin] (store) {7. Αποθήκευση εγκεκριμένης εγγραφής σε δομημένο JSON};
\node[main, below=0.28cm of store] (app) {8. Αξιοποίηση από διαδικτυακή εφαρμογή και μετρικές αξιολόγησης};

\draw[arrow] (sources) -- (snapshot);
\draw[arrow] (snapshot) -- (relevance);
\draw[arrow] (relevance) -- (extract);
\draw[arrow] (extract) -- (quality);
\draw[arrow] (quality) -- (admin);
\draw[arrow] (admin) -- (store);
\draw[arrow] (store) -- (app);
\end{tikzpicture}

    \end{adjustbox}
    \caption[Μεθοδολογική αρχιτεκτονική της προτεινόμενης προσέγγισης]{Μεθοδολογική αρχιτεκτονική της προτεινόμενης προσέγγισης: επιλογή πηγών, επεξεργασία με LLM, έλεγχος ποιότητας, ανθρώπινη εποπτεία και αξιοποίηση μέσω διαδικτυακής εφαρμογής. Ιδία επεξεργασία.}
    \label{fig:methodology-architecture}
\end{figure}

\subsection{Βήματα της μεθοδολογίας}

Η μεθοδολογική ροή οργανώνεται σε επτά βήματα. Κάθε βήμα περιγράφεται εδώ στο
επίπεδο της μεθόδου. Οι τεχνικές λεπτομέρειες υλοποίησης παρουσιάζονται στο
επόμενο κεφάλαιο.

\textbf{Βήμα 1: Επιλογή πηγών.}
Επιλέγονται διευθύνσεις URL από επίσημες, τραπεζικές ή θεσμικές πηγές που σχετίζονται με
χρηματοδότηση ενεργειακών παρεμβάσεων σε κατοικίες. Η επιλογή μπορεί να γίνει
χειροκίνητα από τον διαχειριστή ή να προκύψει από ήδη γνωστές σχετικές πηγές.

\textbf{Βήμα 2: Συλλογή και καθαρισμός περιεχομένου.}
Το περιεχόμενο της πηγής μετατρέπεται σε κείμενο κατάλληλο για επεξεργασία. Ο
στόχος δεν είναι η πλήρης αναπαραγωγή της ιστοσελίδας, αλλά η απομόνωση του
χρήσιμου κειμένου που περιγράφει το πρόγραμμα ή το προϊόν.

\textbf{Βήμα 3: Απόφαση συνάφειας.}
Η πηγή εξετάζεται ως προς το αν ανήκει στο πεδίο της εργασίας. Ως σχετικές
θεωρούνται πηγές για ενεργειακή αναβάθμιση κατοικιών, οικιακές ΑΠΕ και πράσινα
χρηματοδοτικά προϊόντα με σαφή ενεργειακό κριτήριο.

\textbf{Βήμα 4: Εξαγωγή δομημένων πεδίων.}
Για τις σχετικές πηγές εξάγονται πεδία όπως όνομα προγράμματος, δικαιούχοι,
κριτήρια επιλεξιμότητας, επιλέξιμες παρεμβάσεις, ποσά, επιτόκια, προθεσμίες
και φορέας διαχείρισης. Η έξοδος οργανώνεται σε κοινό σχήμα JSON, ώστε τα
προγράμματα να μπορούν να συγκριθούν και να παρουσιαστούν με ενιαίο τρόπο.
Τα πεδία αυτά ομαδοποιούνται εννοιολογικά σε στοιχεία ταυτότητας, επιλεξιμότητα,
χρηματοδότηση, χρονικά στοιχεία, τεχνικό αντικείμενο και μεταδεδομένα ελέγχου.

\textbf{Βήμα 5: Έλεγχος ποιότητας.}
Η παραγόμενη εγγραφή ελέγχεται ως προς τη δομή, την πληρότητα και τη σύνδεση με
την πηγή. Αν μια πληροφορία δεν υπάρχει στο αρχικό κείμενο, πρέπει να παραμένει
κενή ή να επισημαίνεται ως μη διαθέσιμη, όχι να συμπληρώνεται με εικασία.
Κάθε εξαγωγή αντιμετωπίζεται επίσης ως χρονικό στιγμιότυπο, επειδή οι οδηγοί,
οι προθεσμίες και οι οικονομικοί όροι μπορούν να αλλάξουν.

\textbf{Βήμα 6: Ανθρώπινη εποπτεία.}
Οι εγγραφές που παράγονται αυτόματα δεν θεωρούνται αμέσως τελικές. Ο
διαχειριστής μπορεί να τις ελέγξει, να τις διορθώσει, να τις εγκρίνει ή να τις
απορρίψει πριν γίνουν διαθέσιμες στον τελικό χρήστη.

\textbf{Βήμα 7: Αξιοποίηση στην εφαρμογή.}
Οι εγκεκριμένες εγγραφές χρησιμοποιούνται από τη δημόσια πλευρά της εφαρμογής
για αναζήτηση, παρουσίαση βασικών πεδίων και ερωταπαντήσεις πάνω στα διαθέσιμα
δεδομένα. Με αυτόν τον τρόπο, ο τελικός χρήστης βλέπει οργανωμένη πληροφορία,
ενώ η πηγή παραμένει διαθέσιμη για επιβεβαίωση.

\subsection{Ρόλοι, αξιολόγηση και σύνδεση με την υλοποίηση}

Στη ροή διακρίνονται δύο ρόλοι. Ο διαχειριστής χειρίζεται τις πηγές, ελέγχει
τα παραγόμενα αρχεία JSON και αποφασίζει ποιες εγγραφές εγκρίνονται. Ο τελικός
χρήστης βλέπει μόνο εγκεκριμένα προγράμματα και μπορεί να θέτει ερωτήσεις πάνω
σε αυτά.
Η διάκριση αυτή είναι βασική για να μη συγχέεται η αυτόματη εξαγωγή με επίσημη
ερμηνεία των όρων ενός προγράμματος.

Η αξιολόγηση της μεθοδολογίας βασίζεται σε τέσσερις άξονες: ορθότητα απόφασης
συνάφειας, πληρότητα εξαγόμενων πεδίων, πιστότητα στην αρχική πηγή και συνέπεια
των απαντήσεων της εφαρμογής. Οι μετρικές και τα πειραματικά αποτελέσματα δεν
αναπτύσσονται αναλυτικά εδώ, επειδή αποτελούν αντικείμενο του κεφαλαίου
αποτελεσμάτων.

Πρακτικά, η αξιολόγηση προβλέπει χειροκίνητα ελεγμένο δείγμα πηγών, ώστε να
συγκριθεί η απόφαση συνάφειας του συστήματος με την ανθρώπινη κρίση. Για τις
πηγές που γίνονται δεκτές, τα παραγόμενα πεδία JSON ελέγχονται ως προς την
πληρότητα και ως προς το αν τεκμηριώνονται από το αρχικό περιεχόμενο. Τέλος, οι
απαντήσεις της εφαρμογής εξετάζονται σε σχέση με τα εγκεκριμένα δεδομένα, ώστε
να ελέγχεται ότι δεν εισάγουν πληροφορία εκτός της αποθηκευμένης εγγραφής.

Το κεφάλαιο της υλοποίησης που ακολουθεί εξειδικεύει τα παραπάνω βήματα σε
συγκεκριμένα υποσυστήματα: συλλογή και καθαρισμό ιστοσελίδων, εξαγωγή JSON,
ταξινόμηση συνάφειας, αποθήκευση αποτελεσμάτων, διεπαφή διαχειριστή και
δημόσια διεπαφή χρήστη. Έτσι, το παρόν κεφάλαιο παραμένει σε επίπεδο μεθόδου,
ενώ οι τεχνικές επιλογές παρουσιάζονται εκεί όπου ανήκουν.
